\documentclass[10pt,twocolumn]{ctexart}

\usepackage[letterpaper,margin=0.72in]{geometry}
\usepackage{graphicx}
\usepackage{booktabs}
\usepackage{array}
\usepackage{url}
\usepackage{xcolor}
\usepackage{caption}
\usepackage{enumitem}

\setlength{\columnsep}{0.22in}
\setlength{\parskip}{2pt}
\setlength{\parindent}{1em}
\setlist[itemize]{leftmargin=*,itemsep=1pt,topsep=2pt}
\setlist[enumerate]{leftmargin=*,itemsep=1pt,topsep=2pt}
\captionsetup{font=small}
\newcommand{\stackline}[2]{\begin{quote}\small\ttfamily #1\\#2\end{quote}}

\title{\vspace{-1.2em}面向 AI Agent 系统效应的语义火焰图}
\author{AgentSight 研究原型}
\date{2026 年 6 月 14 日}

\begin{document}
\maketitle
\vspace{-1.4em}

\begin{abstract}
AI 编程 agent 的困难不只是 ``它说了什么''，而是它把用户意图转化成了哪些真实系统效应：进程、文件、网络、测试、失败重试和 token 消耗。
现有 LLM observability 工具擅长展示单次运行的 trace tree，传统 profiler 擅长展示 CPU 或内存热点，但二者都难以回答一个用户更常问的问题：哪些任务语义导致了最多、最绕或最重复的系统行为？
本文提出一种语义火焰图模型：用本地小模型给 session、prompt 和 LLM call 生成一个词标签；让 tool/process/file/network 足迹通过确定性 lineage 或 agent-native tool log 继承标签；再把路径折叠成 flamegraph stacks。
在本机一个 AgentSight 仓库的 36 个 Codex/Claude 真实 session 上，agent-native 原型将 4031 个 tool event 展开为 5312 个系统观测并折叠成 2270 条语义栈，压缩 2.34 倍；去掉 session/prompt 语义后，392 个 nonsemantic buckets 会混合多个语义区域，覆盖 68.505\% 的观测权重，flat effect baseline 的混合权重达到 74.473\%。
这些结果说明语义帧能在该 workload 中分开普通进程摘要会合并的行为区域。
本文也明确边界：用户效用实验和 live exact AgentSight capture 尚未完成；当前 C6 只在 4 个 effect 的 fixture 上验证 checker 接口。
\end{abstract}

\section{问题：agent 的系统效应不可读}

一个 agent 运行结束后，用户通常不想从第一条消息回放到最后一条消息。
他们想知道：这次运行哪里最重，哪里绕路，哪里反复试错，哪些文件或网络目标被碰过，哪类 prompt 导致了最多系统动作。
传统工具给出的答案都只覆盖一半。
进程工具能告诉我们 \texttt{git}、\texttt{cargo}、\texttt{sed} 或 \texttt{gh} 跑过；LLM trace 工具能告诉我们模型调用、tool call 和 token；普通 flame graph 能把采样栈按宽度聚合。
但是用户真正关心的是跨层关系：``为了修测试，agent 反复运行了哪些命令并读取了哪些路径''，或者 ``为了写论文，Claude-heavy 的 \texttt{gh} 行为集中在哪些 prompt 区域''。

本文的立场是：agent observability 的第一对象不应该是单个 span，也不应该是单个进程，而应该是 \emph{task-grounded system effect}。
一个面向审计的图必须同时满足三点：

\begin{itemize}
  \item 上层可读：分析者能看到 session/prompt/subagent/LLM call 对应的任务语义。
  \item 下层可证：exact stream 中的系统效应应能追溯到 tool call、进程、文件或网络事件；agent-native prototype 至少保留 tool/log 足迹。
  \item 横向可聚合：重复行为必须折叠成带权重的 stack，而不是画成越来越长的 trace tree。
\end{itemize}

这个目标与现有 LLM observability 平台互补。
LangSmith、Langfuse、Phoenix、Weave、Helicone 等工具已经证明了 LLM trace 的价值：它们记录模型调用、tool calls、retrieval、latency、cost 和 nested spans。
SigNoz/Inkeep 甚至已经把 agent spans 画成 flamegraph，用宽度表示 duration。
这些系统解决的是 ``单次运行如何发生''。
本文试图解决的是另一个问题：\emph{跨 session、跨 agent、跨 prompt，哪些语义区域造成了哪些系统足迹}。

\section{核心模型}

\begin{figure}[t]
  \centering
  \includegraphics[width=\linewidth]{figures/fig-model.pdf}
  \caption{语义标签只命名上层上下文；系统足迹不由 LLM 猜测，而是通过 exact lineage 或 agent-native tool log 继承标签并折叠成 stack。}
  \label{fig:model}
\end{figure}

图~\ref{fig:model} 展示了本文的模型。
语义层非常小：每个 session、prompt 和 LLM call 只生成一个 lowercase word，例如 \texttt{paper}、\texttt{debug}、\texttt{flamegraph}。
这些词不是 ontology，也不是固定分类集合；它们只是给折叠栈提供一个短、稳定、可搜索的任务帧。

系统层不交给模型判断。
对于 agent-native 历史，原型从 tool call 中抽取 tool 类别、命令 basename、effect 类别、路径或域名组。
对于 AgentSight exact stream，目标输入是：

\begin{quote}\small
\texttt{tool\_call -> shell -> child process -> file/network effect}
\end{quote}

每个下层事件必须通过 tool call id、进程实例、child-process family 和时间窗口连接到 prompt/session。
PID-only match 不足够，因为 live tracing 中 PID 会复用。
因此系统效应继承语义标签，而不是由 LLM 对文件或网络事件重新分类。

最终 system footprint stack 为：

\stackline{project;agent;session;prompt;tool;cmd}{effect;path/domain/status}

token footprint stack 为：

\stackline{project;agent;session;prompt;llm;model}{kind}

火焰图的宽度是折叠后的权重。
如果同一路径出现 75 次，folded file 中只有一行带权重的 stack。
这是本文与 trace tree 的根本区别：trace tree 保留时序展开，semantic flamegraph 保留可聚合的 task-effect 结构。

\subsection{为什么只要一个词}

本文刻意不让小模型输出长摘要、层级分类或自然语言解释。
原因有三点。
第一，火焰图帧必须短；长标签会让最宽区域不可读，也会把视觉空间浪费在语言表述差异上。
第二，标签只服务于定位和聚合，不承担判决功能。
如果 \texttt{paper} 与 \texttt{writing} 被分成两个词，最坏结果是相邻区域没有合并；但如果让模型判断某次文件写入 ``是否合理''，错误会直接变成安全或正确性结论。
第三，一个词标签可以本地化、缓存化和可替换。
在当前原型中，小模型只处理 session/prompt/LLM-call 文本一次，系统事件完全通过 lineage 继承标签。
这使得昂贵或隐私敏感的 prompt 处理停留在语义边界内，后续 file/network/process 聚合可以用普通程序完成。

这种设计也给论文 claim 设了边界。
语义标签不是 ground truth task ontology，而是 lossy navigation frame。
因此 evaluation 不问 ``标签是否代表真实意图''，而先问一个更可验证的问题：加入这些帧后，普通工具会合并掉的系统行为是否被分开。
这正是 baseline mixing 指标要测量的内容。

\section{实现}

原型位于 \texttt{docs/visexp}。
一键入口为：

\begin{quote}\small
\texttt{python3 docs/visexp/run\_pipeline.py --out docs/visexp/out}
\end{quote}

该 pipeline 读取本机 Codex 与 Claude JSONL session，运行语义标注、folded stack 生成、SVG/论文图生成、tag stability smoke、C5 task packet 生成、C6 fixture lineage checker、artifact evaluator、visual summary、verifier 和 unit tests。
在当前机器上，完整 pipeline 8 步耗时 86.606 秒，其中 semantic flamegraphs 36.657 秒，tag stability 49.702 秒。

\subsection{小模型标签}

标签器优先调用 \texttt{llama.cpp}。
默认发现到 \texttt{qwen2.5-3b-instruct-q4\_k\_m.gguf} 时，对最多前 60 个 uncached 请求使用本地小模型，temperature 为 0，并用 grammar 约束输出一个词。
非法输出会被丢弃并回退到 deterministic regex tagger。
为了避免 fallback 缓存污染模型结果，cache key 包含 tagger namespace，即 \texttt{fallback} 或模型 basename。

当前运行中共有 3243 个 tag 请求，其中 3213 个命中缓存，30 个本地小模型 uncached 调用全部成功，fallback 使用为 0。
这证明小模型路径是可运行的，但不能证明语义充分；C2 只支持 one-word grammar 和 provenance，C7 的 adequacy 仍需要人工标签。

\subsection{多维投影}

除了完整 system/token flamegraph，原型还输出五个投影：
\texttt{session-system}、\texttt{prompt-system}、\texttt{session-token}、\texttt{prompt-token}、\texttt{llm-token}。
投影不改变总权重，只删除一个或多个语义轴。
例如 \texttt{session-system} 删除 prompt 帧，用来观察同一 session 下系统行为是否跨 prompt 聚合；
\texttt{prompt-system} 删除 session 帧，用来观察同名 prompt tag 是否在不同 session 中共享系统足迹。
verifier 强制检查这些投影的总权重必须与原始 system/token 总量一致。

\subsection{隐私与复现}

提交 artifact 不包含原始 prompt。
\texttt{prompt-tags.csv} 中 prompt preview 默认 redacted，只保留 hash、tag 和计数。
\texttt{verify\_artifacts.py} 检查 prompt redaction、folded totals、manifest hash、dimension projection totals 和敏感字符串。
\texttt{input-manifest.json} 记录 argv、session content hash、脚本 hash、模型 checksum 和 repo dirty 状态。

\subsection{设计取舍}

实现中最重要的取舍是把 ``解释'' 与 ``关联'' 分开。
LLM 只解释上层文本上下文；确定性代码负责解析工具调用、展开命令、折叠路径、生成火焰图和校验总量。
这避免了一个常见陷阱：让 LLM 读完整 trace 后输出一个看似合理的故事。
那样的故事很容易吸引人，却无法证明每个文件、网络或测试动作来自哪次 tool call。
本文需要的是可审计 profile，因此每条 folded stack 都必须能回到一个输入事件集合。

第二个取舍是保留多个视图而不是寻找唯一正确图。
\texttt{session-system} 适合做高层审计，\texttt{prompt-system} 适合定位导致重系统行为的 prompt 类型，\texttt{llm-token} 适合看模型调用的 token accounting。
这些图共享同一套 folded-stack grammar 和 verifier，所以不是三个互不相干的 dashboard，而是同一事实表上的不同投影。
如果某个投影的总权重与源图不一致，verifier 会失败。
这让可视化从 ``漂亮截图'' 变成可测试 artifact。

第三个取舍是承认 agent-native history 和 exact system stream 的差距。
Codex/Claude JSONL 能较好恢复 tool、命令和部分路径域名，但它不是内核级审计日志。
AgentSight 的独特点在于已经能把 \texttt{tool\_call -> shell -> child process -> file/network effect} 精确连起来。
因此本文把 agent-native flamegraph 作为可复现实验输入，把 AgentSight-shaped checker 作为 exact lineage 接口；下一步要用 live snapshot 把二者合并。

\section{评估}

\subsection{评估原则}

我们把 claim 分成三类。
第一类是 artifact-internal claim：folded totals 是否一致、prompt 是否被 redacted、维度投影是否守恒、标签是否满足一个词 grammar。
这类 claim 可以由 verifier 和 unit tests 完成。
第二类是 partitioning claim：语义帧是否分离了 baseline 会混合的行为。
这类 claim 由 nonsemantic mixing 和 flat mixing 度量支撑。
第三类是 user/system claim：用户是否更快回答问题，以及 live AgentSight capture 是否完整连接 effect lineage。
这类 claim 目前不能靠本地 artifact 自证，所以在本文中必须保守呈现。

这种 gate 化写法使得负结果不会污染正结果。
即使 C5 用户实验尚未完成，C1--C3 仍然可以成立：它们只依赖本机真实 session、折叠栈守恒和 baseline 对比。
反过来，如果后续用户实验不显著，semantic flamegraph 仍可能是一个好的 forensic profile，只是不能宣称提升普通用户效率。

\begin{figure}[t]
  \centering
  \includegraphics[width=\linewidth]{figures/fig-results.pdf}
  \caption{当前 artifact 的主要结果。Collapsed 表示语义 system stack 的折叠观测占比；Nonsem. mix 和 Flat mix 表示去掉语义帧后 baseline bucket 混合多个 prompt/session 的观测权重占比。}
  \label{fig:results}
\end{figure}

\begin{figure}[t]
  \centering
  \includegraphics[width=\linewidth]{figures/fig-dimensions.pdf}
  \caption{不同语义维度投影的压缩率。Token 视图因为按 token 权重聚合，压缩率远高于 system event 视图；这些数字只能作为 source-local accounting。}
  \label{fig:dimensions}
\end{figure}

\subsection{RQ1：语义帧是否分割 baseline 混合？}

表~\ref{tab:metrics} 和图~\ref{fig:results} 给出当前本机结果。
36 个 session 包含 4031 个 raw tool event。
路径和域名展开后得到 5312 个 system observations，折叠成 2270 条 unique semantic stacks，压缩率为 2.34 倍。
更关键的是 baseline mixing：如果去掉 session/prompt 帧，392 个 nonsemantic buckets 会混合多个语义区域，覆盖 68.505\% 的观测权重；如果进一步退化成 flat effect bucket，397 个 buckets 混合多个语义区域，覆盖 74.473\% 的观测权重。
这个指标是 proxy：它证明语义帧在该 workload 中分割了 baseline 会合并的 buckets，但不证明用户效用或标签语义正确。

\begin{table}[t]
  \centering
  \small
  \caption{当前 artifact 指标。}
  \label{tab:metrics}
  \begin{tabular}{lr}
    \toprule
    指标 & 数值 \\
    \midrule
    Sessions & 36 \\
    Raw tool events & 4031 \\
    Expanded system observations & 5312 \\
    Unique semantic system stacks & 2270 \\
    Collapsed observations & 3042 \\
    Semantic compression & 2.34x \\
    Unique prompt tags & 38 \\
    Generic prompt rows & 35.987\% \\
    Nonsemantic mixed weight & 68.505\% \\
    Flat mixed weight & 74.473\% \\
    \bottomrule
  \end{tabular}
\end{table}

这个结果直接回应 ``传统工具也能看到'' 的质疑，但只在 artifact-level 成立。
传统 flat summary 可以告诉我们 \texttt{git read} 或 \texttt{cargo test} 很重，但不能告诉这些重行为分别来自 \texttt{paper}、\texttt{debug}、\texttt{flamegraph} 还是 \texttt{session} prompt 区域。
混合 bucket 的高权重说明语义帧不是装饰，而是在该 workload 的 baseline 中被合并掉的信息。

\subsection{Case study：本仓库 session 能看出什么}

表~\ref{tab:case} 展示了当前运行中最宽的一组真实 stack。
这些 stack 不包含原始 prompt，却保留了足够的任务上下文。
例如，传统 process summary 会说 Codex 的 \texttt{git read} 出现 376 次；语义 stack 进一步显示，这个 bucket 实际混合了 27 个 session/prompt 区域。
最大几个区域分别是 \texttt{session/session}、\texttt{flamegraph/design}、\texttt{paper/debug}、\texttt{flamegraph/branch} 和 \texttt{session/trace}。
另一个例子是 \texttt{cargo test}：82 次测试行为被 baseline 合并，但语义帧把它分成 session 解析、paper debug、model 检查和 SSL/debug 等区域。

\begin{table}[t]
  \centering
  \small
  \caption{真实本机 session 的 top semantic stacks 摘要。}
  \label{tab:case}
  \begin{tabular}{@{}p{0.34\linewidth}p{0.48\linewidth}r@{}}
    \toprule
    语义区域 & 系统行为 & 权重 \\
    \midrule
    Codex \texttt{session}/\texttt{session} &
    \texttt{git} read, status ok & 75 \\
    Codex \texttt{flamegraph}/\texttt{design} &
    \texttt{git} read, status ok & 56 \\
    Codex \texttt{paper}/\texttt{debug} &
    \texttt{git} read, status ok & 38 \\
    Codex \texttt{session}/\texttt{session} &
    \texttt{sed} read \texttt{collector/src/sources} & 34 \\
    Codex \texttt{flamegraph}/\texttt{design} &
    \texttt{python3} on semantic flamegraph tests & 32 \\
    \bottomrule
  \end{tabular}
\end{table}

这个例子也说明本文不能只画一棵 trace tree。
trace tree 能重放某次 \texttt{git} 调用的位置，但不能自然回答 ``跨多个 prompt，哪些 \texttt{git read} 属于同一调试或设计阶段''。
flat process summary 可以聚合 \texttt{git read}，但会丢掉任务区域。
semantic flamegraph 把二者连起来：宽度来自重复系统行为，层级来自 session/prompt/tag lineage。
因此它适合做三类 artifact audit：重区域，例如 \texttt{session/session} 的 \texttt{git read} 行为；绕路区域，例如同一 prompt tag 下反复测试或反复读取相同路径；重复区域，例如多个 session 都把相同 collector 路径作为读取热点。
这些是审计 agent 工作量时会问的问题，而不是单纯的性能热点。

\subsection{RQ2：单独看 session/prompt/LLM 维度有用吗？}

图~\ref{fig:dimensions} 展示了维度投影。
\texttt{session-system} 有 1603 条 unique stacks，压缩 3.314 倍；
\texttt{prompt-system} 有 2170 条 unique stacks，压缩 2.448 倍。
这说明 session 是更粗的聚合单位，适合回答 ``这次会话整体做了什么''；prompt 是更细的聚合单位，适合定位 ``哪类用户请求导致了某个系统足迹''。

Token 维度压缩率更高：
\texttt{session-token} 压缩 5293023.833 倍，\texttt{llm-token} 压缩 1456795.55 倍。
这不是性能结论，而是说明 token 是天然大权重聚合量。
因此论文只把 token flamegraph 用作 source-local accounting，不做跨 agent 成本比较。

\begin{figure}[t]
  \centering
  \includegraphics[width=\linewidth]{figures/fig-flame-excerpt.pdf}
  \caption{系统语义火焰图的 top-stack 摘要。完整交互式候选图在 \texttt{docs/visexp/out/system-flamegraph.svg}。}
  \label{fig:flame}
\end{figure}

\subsection{RQ3：标签是否稳定？}

当前 C7 只能给 partial verdict。
语法上，prompt tag invalid count 为 0，same-hash prompt tag conflict 为 0。
tag stability smoke 对 24 个 session/prompt/LLM fragments 重复 2 次，fallback 与 llama 路径在本次 run 中都没有 invalid output。
但是 generic prompt row share 仍有 35.987\%。
这意味着 one-word tags 现在足够作为导航 hint，但还不足以作为语义 ontology。
顶会版本需要人工 adequacy labels，并报告模型间一致性和 generic-rate 降低策略。

\subsection{RQ4：exact lineage 接口是否可执行？}

作为接口烟测，C6 的 checker 已经可执行。
在 AgentSight-shaped fixture 中，4 个 process/file/network effects 全部连接到 session/tool/prompt ancestry，join rate 为 100.0\%，并输出 exact-effect folded stacks。
但这只是 fixture-only checker evidence，不属于主 workload 指标。
本文不能声称 live exact capture 已完成；真正的 C6 需要对真实 AgentSight session snapshot 运行同一 checker，并把 unjoined in-scope effect 作为 collector 或 join bug。

\section{相关工作}

\textbf{Flame graphs.}
Brendan Gregg 的 flame graph 把采样栈折叠成宽度可读的性能图，核心思想是 ``找最宽的块''。
Grafana、Datadog、Splunk 等系统也把 profiling flamegraph 用于 CPU、内存或 profile drilldown。
本文借用 folded stack 形式，但 stack frame 不来自函数调用栈，而来自 agent semantic context 和 system-effect lineage。

\textbf{Distributed tracing.}
Dapper 与 OpenTelemetry 体系将请求路径表示为 spans。
LLM/agent observability 平台也普遍采用 trace/span tree。
LangSmith 的 runs、Langfuse 的 nested observations、Phoenix 的 OpenTelemetry/OpenInference traces、Weave 的 Calls/Traces/Threads 都能帮助开发者调试单次请求。
本文的差异是把 trace 转成可聚合的 task-effect profile，并用 baseline mixing 衡量语义帧能分割哪些 baseline buckets。

\textbf{Agent observability.}
Helicone、Braintrust、Weave、Phoenix、LangSmith 等系统都支持 tool-call tracing、session grouping、cost/latency 和 evaluation。
SigNoz/Inkeep 的 agent flamegraph 已经把多 agent workflow spans 以 duration 宽度展示，适合发现 critical path、error cascade 和 tool overhead。
本文与它们的互补点在于：我们不只画单次运行的 span duration，而是跨 session 聚合 prompt/session/LLM tags 与系统效应，回答 ``哪些任务语义造成系统足迹''。

\textbf{Context and token profiling.}
ContextFlame、context-profiler 等工作关注 LLM context window 和 token hotspot。
本文的 token flamegraph 与其方向相近，但系统贡献在于把 token 视图与 tool/process/file/network effect 视图放入同一个语义栈模型。

\section{局限与下一步}

当前版本离 OSDI 强接收还有三个缺口。
第一，C5 用户效用没有真实参与者结果。
我们已经生成 6 个任务、answer key、participant packets、response template 和 scorer，但没有 participant CSV。
第二，C6 没有 live exact AgentSight capture。
pipeline 已支持 \texttt{--lineage-snapshot}，但当前 artifact 只提交 fixture checker。
第三，C7 标签质量还不够。
35.987\% generic prompt rows 会削弱导航价值，必须通过更好的小模型、batch annotation、label repair 或人工 adequacy gate 降低。

这些局限也定义了下一轮实验。
如果 4 人 pilot 发现 semantic flamegraph 在 repeated/heavy/divergent tasks 上提高准确率或降低时间，就可以把 C5 从 unsupported 提升到 partial。
如果 3 个 live AgentSight sessions 的 in-scope effects 全部 join 到 prompt/tool ancestry，C6 才能成为论文的系统贡献。
如果人工 adequacy labels 显示 one-word tags 足以导航但不等于 ontology，论文应把标签定位为 lossy navigation layer，而不是语义真值。

\section{结论}

本文提出 semantic flamegraph：用本地小模型给 agent session/prompt/LLM call 生成短语义帧，用 lineage 或 agent-native tool log 连接系统足迹，再用 folded stack 聚合跨 session 行为。
当前本机 artifact 显示，语义帧能分离传统 process summary 合并掉的大量行为区域：nonsemantic baseline 的 mixed weight 为 68.505\%，flat baseline 为 74.473\%。
因此 novelty 不在 ``又画了一张 flamegraph''，而在把任务语义、LLM/tool history 与可测试的系统足迹投影纳入同一个可折叠模型。
当前证据足以支持 artifact-level partitioning claim，不足以支持用户效用或 live exact capture 的强 claim。

\begin{thebibliography}{9}
\bibitem{gregg-flame}
Brendan Gregg.
Flame Graphs.
\url{https://www.brendangregg.com/flamegraphs.html}.

\bibitem{gregg-ai}
Brendan Gregg.
AI Flame Graphs.
\url{https://www.brendangregg.com/blog/2024-10-29/ai-flame-graphs.html}.

\bibitem{dapper}
Benjamin H. Sigelman et al.
Dapper, a Large-Scale Distributed Systems Tracing Infrastructure.
Google Technical Report, 2010.

\bibitem{otel-genai}
OpenTelemetry.
Semantic Conventions for GenAI agent and framework spans.
\url{https://github.com/open-telemetry/semantic-conventions-genai}.

\bibitem{langsmith}
LangChain.
LangSmith Observability Concepts.
\url{https://docs.langchain.com/langsmith/observability-concepts}.

\bibitem{langfuse}
Langfuse.
Observability data model.
\url{https://langfuse.com/docs/observability/data-model}.

\bibitem{phoenix}
Arize Phoenix.
What is Phoenix?
\url{https://arize.com/docs/phoenix}.

\bibitem{weave}
Weights \& Biases.
What is Weave?
\url{https://docs.wandb.ai/weave/concepts/what-is-weave}.

\bibitem{signoz-inkeep}
SigNoz.
How Inkeep Monitors Their AI Agent Framework with SigNoz.
\url{https://signoz.io/blog/inkeep-ai-agent-monitoring/}.
\end{thebibliography}

\end{document}
